% !TeX program = tectonic
\documentclass[11pt,a4paper]{article}
\usepackage{fontspec}
\usepackage[russian]{babel}
\babelfont{rm}[Language=Default]{Liberation Serif}
\babelfont[Language=Default]{sf}{Carlito}
\babelfont[Language=Default]{tt}{DejaVu Sans Mono}
\usepackage{amsmath,amssymb,amsthm}
\usepackage{geometry}
\geometry{a4paper, top=2.4cm, bottom=2.4cm, left=2.4cm, right=2.4cm}
\usepackage{booktabs,tabularx,enumitem,xcolor,tcolorbox}
\tcbuselibrary{breakable,skins}
\definecolor{linkblue}{HTML}{1F4E79}
\definecolor{statgreen}{HTML}{2F6B3A}
\definecolor{goldacc}{HTML}{8B7E5A}
\usepackage[colorlinks=true,linkcolor=linkblue,urlcolor=linkblue,unicode=true]{hyperref}
\theoremstyle{plain}
\newtheorem{theorem}{Теорема}
\newtheorem{lemma}[theorem]{Лемма}
\newtheorem{proposition}[theorem]{Предложение}
\newtheorem{corollary}[theorem]{Следствие}
\theoremstyle{definition}
\newtheorem{definition}[theorem]{Определение}
\theoremstyle{remark}
\newtheorem{remark}[theorem]{Замечание}
\newtcolorbox{statusbox}[1][]{colback=statgreen!5,colframe=statgreen!75,
  title={\textbf{Сводка доказанного:} #1},fonttitle=\small,boxrule=0.7pt,arc=2pt,breakable}
\newtcolorbox{databox}[1][]{colback=goldacc!6,colframe=goldacc!80,
  title={\textbf{Данные протокола:} #1},fonttitle=\small,boxrule=0.7pt,arc=2pt,breakable}
\newtcolorbox{verifybox}[1][]{colback=linkblue!5,colframe=linkblue!80,
  title={\textbf{Верификация:} #1},fonttitle=\small,boxrule=0.7pt,arc=2pt,breakable}
\widowpenalty=10000
\clubpenalty=10000
\setlength{\parskip}{0.35em}
\newcommand{\CC}{\mathbb{C}}
\newcommand{\RR}{\mathbb{R}}
\newcommand{\ZZ}{\mathbb{Z}}
\newcommand{\QQ}{\mathbb{Q}}
\newcommand{\Ga}{\Gamma}
\newcommand{\Bfun}{\mathrm{B}}
\newcommand{\im}{\operatorname{Im}}
\newcommand{\re}{\operatorname{Re}}
\newcommand{\lcm}{\operatorname{lcm}}
\newcommand{\SNF}{\operatorname{SNF}}
\newcommand{\Jac}{\operatorname{Jac}}
\newcommand{\rank}{\operatorname{rank}}
\newcommand{\Ind}{\operatorname{Index}}
\newcommand{\disc}{\operatorname{disc}}
\begin{document}
\begin{center}
{\LARGE\bfseries Замкнутая форма периода}\\[0.4em]
{\large бета-интеграл, фазовая решётка и эквивариантность}\\[0.6em]
{\normalsize Исаев Исхак Хамзатович}\\[0.2em]
{\small Программа <<Динамический принцип>> --- лаборатория \texttt{hodge-laboratory}}\\[0.15em]
{\small Отдельная монография теоремы · Монография 2/16}
\end{center}
\vspace{0.4em}
{\color{linkblue}\hrule height 0.8pt}
\vspace{0.8em}

\section{Постановка}

Периоды кривой Ферма --- интегралы голоморфных форм по циклам ---
допускают замкнутую форму: трансцендентная часть (бета-интеграл) и
алгебраическая часть (корень единицы) разделяются явно. Это
разделение --- техническое основание всей $\Ga$-нормировки
программы: не имея замкнутой формы, нормировка была бы численной
подгонкой.

\begin{lemma}[базис голоморфных форм]\label{lem:forms}
Формы
\[
  \eta_{i,j} \;=\; \frac{x^i\,y^j\,dx}{y^{\,N-1}},
  \qquad i,j\ge 0,\quad i+j\le N-3,
\]
голоморфны на $C_N$; их число равно $g$, и они образуют базис
$H^0(C_N,\Omega^1)$. Каждая $\eta_{i,j}$ --- собственная форма
действия $\mu_N\times\mu_N$ с характером $(a,b)=(i+1,j+1)$.
\end{lemma}

\begin{lemma}[замкнутая форма периода]\label{lem:closedform}
В координате $t=x^N$ справедливо
\[
  \eta_{i,j} \;=\; \frac{1}{N}\,
  t^{\frac{a}{N}-1}(1-t)^{\frac{b}{N}-1}\,dt,
  \qquad a=i+1,\ b=j+1.
\]
Для подъёма $\gamma$ пути с данными обхода $(r,s)\in(\ZZ/N)^2$
период имеет замкнутую форму
\[
  P_{a,b}(r,s) \;=\; \frac{1}{N}\,\zeta_N^{\,ra+sb}\;\Omega_{a,b},
  \qquad
  \Omega_{a,b}\;=\;\frac{\Ga(\frac aN)\Ga(\frac bN)}{\Ga(\frac{a+b}N)}
  \;>0.
\]
Образ период-функционала лежит в решётке
$\frac{\Omega_{a,b}}{N}\ZZ[\zeta_N]$.
\end{lemma}

\begin{proof}
Подстановка $t=x^N$ даёт $dt=Nx^{N-1}dx$, значит
$x^i\,dx=\tfrac1Nt^{\frac{i+1}N-1}dt$; и
$y^j/y^{N-1}=(1-t)^{\frac{j+1}N-1}$, так как $y^N=1-t$. Каждая
полная петля вокруг $t=0$ умножает ветвь $t^{1/N}$ на $\zeta_N$ ---
интеграл умножается на $\zeta_N^a$; петля вокруг $t=1$ даёт
$\zeta_N^b$. Путь с данными обхода $(r,s)$ несёт множитель
$\zeta_N^{ra+sb}$ на эталонном интеграле
$\frac1N\int_0^1t^{\frac aN-1}(1-t)^{\frac bN-1}dt
=\frac1N\Bfun(\frac aN,\frac bN)$. Голоморфность: в точке
$(\zeta,0)$ имеем $dx\sim-\zeta y^{N-1}dy$, форма регулярна; в
бесконечности в координатах $u=X/Y$, $v=Z/Y$ форма равна
$u^i(v\,du-u\,dv)v^{N-3-i-j}$ --- регулярна при $i+j\le N-3$.
Независимость: формы с различными характерами независимы
(проекционные операторы по группе). Решётка: любой цикл ---
целочисленная комбинация замкнутых цепей, интеграл по ребру ---
разность двух значений замкнутой формы.
\end{proof}

\begin{corollary}[$\mu_N\times\mu_N$-эквивариантность]
Перенос цикла действием $(\zeta^u,\zeta^v)$ умножает период
собственной формы $(a,b)$ на $\zeta_N^{ua+vb}$.
\end{corollary}

\begin{databox}[числа]
$N=15$, $(a,b)=(1,1)$: $\Omega_{1,1}=\Ga(1/15)^2/\Ga(2/15)
=246.897\ldots$; $N=15$, $(a,b)=(7,8)$: $\Omega_{7,8}
=\pi/\sin84^\circ=3.16298\ldots$ --- чистая дробная доля $\pi$.
Сертификат B (V2): максимум относительного отклонения замкнутой
формы от независимого tanh-sinh на $69$ тестах ---
$3.9\cdot10^{-36}$; фазы (V3) --- отклонение $0.0$ радиан.
\end{databox}

\begin{statusbox}[итог]
Замкнутая форма периода доказана подстановкой и анализом ветвлений;
эквивариантность --- прямое следствие. Фаза фиксируется явно ---
никакая <<общая фаза>> не ищется и не постулируется.
\end{statusbox}

\begin{verifybox}[как воспроизвести]
\texttt{python3 laboratory.py} $\to$ <<Стенд N=15/30 $\to$
V2/V3/V4>>; конструктор экспериментов: произвольные $N$, $(a,b)$,
$(r,s)$ --- лаборатория печатает замкнутую форму, независимый
интеграл и вердикт.
\end{verifybox}

\vfill
{\color{linkblue}\hrule height 0.6pt}
\vspace{0.5em}
{\small\color{gray}
Лицензия: индивидуальная исключительная лицензия Исаева Исхака Хамзатовича.
Все права на вычисления, формулы и тексты принадлежат автору.
Верификация: \texttt{python3 laboratory.py} (пункт <<Стенды>>/<<Сертификаты>>),
Lean: \texttt{verification/lean/HodgeLaboratory.lean}.
}
\end{document}
